% =====================================================================
% AI Academy -- National AI Center
% Software Engineering Final Project -- Team Report
% Topic 4: Async Research Assistant
% =====================================================================

\documentclass[11pt,a4paper]{article}

\usepackage[dvipsnames]{xcolor}
\definecolor{accent2}{HTML}{8B0000}
\definecolor{ph}{HTML}{6B7280}
\newcommand{\phh}[1]{\textcolor{ph}{\textit{#1}}}

% ===== CONFIGURATION ====================================
\newcommand{\teamname}{Individual Submission}
\newcommand{\topicchoice}{Topic 4 --- Async Research Assistant}
\newcommand{\member}[2]{\textbf{#1} (#2)}
\newcommand{\reportdate}{May 23, 2026}
\newcommand{\repolink}{github.com/emiljafarov3841/async-research-assistant}

\newcommand{\coursecode}{AI-ENG-110}
\newcommand{\coursename}{Software Engineering}
\newcommand{\institution}{AI Academy, National AI Center}
\newcommand{\semester}{Spring 2026}

% ===== PACKAGES ======================================================
\usepackage[margin=2.2cm,headheight=15pt]{geometry}
\usepackage[T1]{fontenc}
\usepackage[utf8]{inputenc}
\usepackage[final,protrusion=true,expansion=false]{microtype}
\usepackage{amsmath,amssymb}
\usepackage{graphicx}
\usepackage{booktabs}
\usepackage{array}
\usepackage{tabularx}
\usepackage{ragged2e}
\usepackage{enumitem}
\usepackage[parfill]{parskip}
\usepackage{listings}
\usepackage[most]{tcolorbox}
\usepackage{titlesec}
\usepackage{fancyhdr}
\usepackage{lastpage}
\usepackage[hidelinks]{hyperref}

\emergencystretch=3em
\hbadness=10000
\hfuzz=2pt

\hypersetup{
  colorlinks=true,
  linkcolor=NavyBlue,
  urlcolor=NavyBlue,
  citecolor=NavyBlue,
  pdftitle={\coursecode\ Final Project Report},
  pdfauthor={\teamname}
}

\definecolor{accent}{HTML}{0B3D91}
\definecolor{soft}{HTML}{F2F4F8}
\definecolor{deliver}{HTML}{E8F0FE}
\definecolor{deliverframe}{HTML}{1A56DB}

\titleformat{\section}
  {\Large\bfseries\color{accent}}{\thesection.}{0.6em}{}
\titleformat{\subsection}
  {\large\bfseries\color{accent}}{\thesubsection}{0.5em}{}

\let\ph\phh

\definecolor{codegreen}{rgb}{0,0.55,0}
\definecolor{codegray}{rgb}{0.4,0.4,0.4}
\definecolor{codepurple}{rgb}{0.55,0,0.78}
\definecolor{backcolour}{rgb}{0.97,0.97,0.97}

\lstdefinestyle{python}{
  language=Python,
  backgroundcolor=\color{backcolour},
  commentstyle=\color{codegreen},
  keywordstyle=\color{accent},
  numberstyle=\tiny\color{codegray},
  stringstyle=\color{codepurple},
  basicstyle=\ttfamily\footnotesize,
  breakatwhitespace=false,
  breaklines=true,
  keepspaces=true,
  numbers=left,
  numbersep=6pt,
  tabsize=2,
  frame=single,
  framerule=0.4pt,
  rulecolor=\color{black!30},
}

\pagestyle{fancy}
\fancyhf{}
\fancyhead[L]{\small\textit{\coursecode\ -- Final Report -- \teamname}}
\fancyhead[R]{\small\textit{\semester}}
\fancyfoot[L]{\small\textit{\institution}}
\fancyfoot[R]{\small Page \thepage\ of \pageref{LastPage}}
\renewcommand{\headrulewidth}{0.4pt}
\renewcommand{\footrulewidth}{0.4pt}

% =====================================================================
\begin{document}

\begin{center}
{\small \textsc{\institution} \\[0.2em] \textsc{\coursecode\ -- \coursename\ -- \semester}}\\[1.2em]
{\Large\bfseries\color{accent} Final Project Report}\\[0.3em]
{\LARGE\bfseries \topicchoice}\\[0.6em]
\rule{0.85\textwidth}{0.6pt}
\end{center}

\vspace{0.6em}
\begin{tabularx}{\textwidth}{@{}lXlX@{}}
\textbf{Team:} & \teamname & \textbf{Date:} & \reportdate \\
\textbf{Repo:} & {\footnotesize \repolink} & \textbf{Tag:} & \texttt{v1.0-final} \\
\textbf{Submitter:} & \multicolumn{3}{X@{}}{\member{Emil Jafarov}{@emiljafarov3841}}\\
\end{tabularx}
\vspace{0.4em}\hrule\vspace{0.6em}

% ===== EXECUTIVE SUMMARY =============================================
\section*{Executive Summary}
\addcontentsline{toc}{section}{Executive Summary}

We built an Async Research Assistant that takes a user question and queries Wikipedia, arXiv, and a web-search API \textbf{in parallel} through a semaphore-bounded asyncio pipeline, then uses a large language model to synthesize a single cited answer. Our headline concurrency result is a \textbf{3.0$\times$ speedup} (sequential 0.767s vs.\ concurrent 0.258s for 5 questions in offline mode with 0.05s simulated latency). Robustness was demonstrated by injecting arXiv timeouts: the system gracefully degraded—continuing with Wikipedia and web results and noting the failure in the output. The single most surprising lesson was that per-source timeout isolation (using \texttt{asyncio.wait\_for()} inside each coroutine rather than a single outer timeout) is what makes graceful degradation actually work in practice.

% ===== TABLE OF CONTENTS =============================================
\tableofcontents
\vspace{0.6em}\hrule\vspace{0.4em}

% =====================================================================
\section{Project Overview}

\subsection{Problem framing \& scope}

The Async Research Assistant addresses the challenge of answering research questions that require synthesizing information from multiple heterogeneous sources. A user submits a question in natural language; the system concurrently retrieves relevant excerpts from Wikipedia (encyclopedic background), arXiv (recent academic papers), and a web search provider (current web content), then asks an LLM to write a concise cited answer.

The CLI entry point is \texttt{python -m researcher ask "your question"}. Users can restrict sources with \texttt{--sources wiki,arxiv} and bypass the cache with \texttt{--no-cache}. A Streamlit web UI (\texttt{streamlit run app.py}) provides a graphical interface as a bonus feature. The system does \textbf{not} provide real-time news or authoritative medical/legal advice, which are known limitations we document explicitly.

We implemented the full spec with one deliberate tightening: the cache backend is a pluggable abstraction (\texttt{CacheBackend} ABC) rather than a single hard-coded implementation, enabling test-time injection of \texttt{InMemoryCache} and production use of \texttt{FilesystemCache}.

\subsection{Provider choice and the AI module contract}

We used \textbf{Anthropic Claude claude-sonnet-4-6} as the default LLM for synthesis and \textbf{Tavily} as the default web search provider for reliable live web results. The provided \texttt{ai/} functions we call are: \texttt{ai.fetch\_wikipedia()}, \texttt{ai.fetch\_arxiv()}, \texttt{ai.fetch\_web()}, and \texttt{ai.synthesize()}. We did \textbf{not} modify the public interface of the provided \texttt{ai/} package.

Defense against malformed responses: the \texttt{AnswerWithCitations} schema from \texttt{ai/schemas.py} already drops out-of-range citation indices in \texttt{\_extract\_cited\_indices}. Our SE layer additionally catches \texttt{ProviderError} and \texttt{ValueError} from synthesis, returning a structured error response rather than a traceback.

% =====================================================================
\section{Architecture}\label{sec:arch}

\subsection{Module map}

The architecture separates the provided \texttt{ai/} boundary from our SE layer:

\begin{center}
\begin{small}
\begin{verbatim}
 CLI / Streamlit UI
       |
       v
 ResearchEngine (src/core/researcher.py)
   |               |
   v               v
 ResearchOrchestrator    ResearchRepository
 (concurrency layer)     (SQLite storage)
   |
   v
 AIService (src/services/ai_service.py)
   |         |
   v         v
 provided    CacheBackend (ABC)
 ai/ package  InMemoryCache | FilesystemCache
\end{verbatim}
\end{small}
\end{center}

The boundary of \texttt{ai/} is crossed only through \texttt{AIService}, which adds retries, timeouts, and logging.

\subsection{OOP design \& key abstractions}

The most important abstraction is the \texttt{CacheBackend} abstract base class:

\begin{lstlisting}[style=python]
class CacheBackend(abc.ABC):
    @abc.abstractmethod
    def get(self, source: str, query: str) -> list[dict] | None: ...

    @abc.abstractmethod
    def set(self, source: str, query: str, data: list[dict]) -> None: ...

    @abc.abstractmethod
    def invalidate(self, source: str, query: str) -> None: ...

    @abc.abstractmethod
    def clear(self) -> None: ...

class InMemoryCache(CacheBackend):
    """Process-local TTL cache backed by a plain dict."""
    ...

class FilesystemCache(CacheBackend):
    """JSON files on disk under CACHE_DIR/{key}.json."""
    ...
\end{lstlisting}

We chose an ABC (not a Protocol) because the two concrete backends share no implementation, and because the ABC enforces at class-definition time that all methods are implemented—Protocol would only catch missing methods at call sites. The \texttt{ResearchOrchestrator} receives a \texttt{CacheBackend} via constructor injection, making it trivially testable by passing an \texttt{InMemoryCache(ttl=0)} in tests.

\subsection{Data flow across module boundaries}

SE-layer models in \texttt{src/models.py}: \texttt{QuestionRequest}, \texttt{ResearchResult}, \texttt{CitationRecord}, \texttt{SourceResult}, \texttt{CacheEntry}, \texttt{ErrorResponse}. No naked dictionaries cross module boundaries. \texttt{ai.Source} objects from the \texttt{ai/} package are immediately converted to \texttt{SourceResult} before being included in \texttt{ResearchResult}.

% =====================================================================
\section{Concurrency \& Performance}\label{sec:concurrency}

\subsection{Why asyncio}

We used \texttt{asyncio} because the bottleneck is I/O (HTTP calls to Wikipedia, arXiv, and the web search API). These are network-bound, not CPU-bound; cooperative multitasking via coroutines has essentially zero overhead compared to thread synchronization. Threads would work but add unnecessary complexity for pure I/O tasks.

The semaphore is set to \texttt{MAX\_CONCURRENT=5} by default (configurable). For a single question, we run exactly 3 concurrent fetches (one per source), so the semaphore primarily prevents unbounded bursts when processing many questions in a batch. Setting the semaphore to 1 would serialize all fetches; setting it to 100 would risk 429s from free-tier APIs.

\subsection{Sequential-vs-concurrent benchmark}

\begin{center}\small
\begin{tabular}{lcccc}
\toprule
\textbf{Workload} & $N$ & \textbf{Sequential (s)} & \textbf{Concurrent (s)} & \textbf{Speedup} \\
\midrule
5 research questions (offline, 0.05s simulated latency) & 5 & 0.767 & 0.258 & 3.0$\times$ \\
5 research questions (live, Anthropic+DDG) & 5 & 8.4 & 3.1 & 2.7$\times$ \\
\bottomrule
\end{tabular}
\end{center}

The speedup is approximately $3\times$ rather than the theoretical maximum because synthesis (the LLM call) is sequential per question and dominates in live mode.

\subsection{Rate limit \& politeness}

We retry transient failures with tenacity (4 attempts, exponential backoff from 1--30s). HTTP 429 and 5xx errors are treated as retryable when they escape directly, and provider-wrapped failures are also retried. In addition, the service layer has a configurable per-source minimum interval (\texttt{SOURCE\_RATE\_LIMIT\_SECONDS}, default 0.2s), so live source calls are not fired in uncontrolled bursts. Tavily is the default web provider because it is more reliable than the no-key DuckDuckGo fallback.

% =====================================================================
\section{Robustness \& Error Handling}\label{sec:robustness}

\subsection{Wrapping the AI module}

Every \texttt{ai.*} call goes through \texttt{AIService}:
\begin{itemize}
  \item \textbf{Retries}: tenacity \texttt{AsyncRetrying}, 4 attempts, exponential backoff 1--30s, retries on \texttt{ConnectionError}, \texttt{TimeoutError}, \texttt{OSError}, \texttt{httpx.TimeoutException}, \texttt{ProviderError}.
  \item \textbf{Timeouts}: per-source \texttt{asyncio.timeout()} wraps each fetcher call (Wikipedia 10s, arXiv 15s, web 10s).
  \item \textbf{Logging}: structured \texttt{logger.info()} at start/end with \texttt{extra=\{"elapsed\_s": ...\}}.
  \item \textbf{Validation}: empty query short-circuits before any network call.
\end{itemize}

\subsection{Failure-mode analysis}

\textbf{Injected failure}: we monkey-patched \texttt{fake\_ai\_service.fetch\_arxiv} to raise \texttt{ConnectionError("arxiv is down")} in \texttt{tests/test\_concurrency.py::test\_orchestrator\_graceful\_degradation}.

\textbf{What the system did}: the bounded \texttt{run\_concurrent(...)} pipeline caught the exception without cancelling the Wikipedia or web fetches. The orchestrator logged \texttt{source\_failed} at WARNING level and added \texttt{"arxiv"} to \texttt{sources\_failed}. The \texttt{ResearchEngine} then proceeded to synthesis with the available sources.

\textbf{What the user saw}: the CLI printed the answer with a note: \textit{"[note] The following sources could not be reached: arxiv"}.

\textbf{What the logs showed}:
\begin{verbatim}
2026-05-14T10:05:21 WARNING  ResearchOrchestrator
  source_failed {"source": "arxiv", "error": "arxiv is down"}
2026-05-14T10:05:22 INFO     ResearchOrchestrator
  orchestrator_done {"total_sources": 2, "failed": ["arxiv"], "elapsed_s": 0.8}
\end{verbatim}

\subsection{Input validation}

\begin{itemize}
  \item \textbf{CLI/API}: \texttt{QuestionRequest} pydantic model enforces \texttt{min\_length=3}, \texttt{max\_length=1000}; invalid \texttt{sources} values raise \texttt{ValueError} before any network call.
  \item \textbf{Empty question}: rejected with \texttt{"question must not be blank"} before touching the orchestrator.
  \item \textbf{Unknown sources}: rejected with \texttt{"Unknown source(s): ['bad']. Allowed: \{arxiv, web, wiki\}"}.
\end{itemize}

% =====================================================================
\section{Storage \& Caching}\label{sec:storage}

\subsection{Schema}

\begin{lstlisting}[style=python,language=sql,numbers=none,basicstyle=\ttfamily\small]
CREATE TABLE IF NOT EXISTS research_sessions (
    id            INTEGER PRIMARY KEY AUTOINCREMENT,
    question      TEXT    NOT NULL,
    answer        TEXT    NOT NULL,
    citations     TEXT    NOT NULL,   -- JSON array of CitationRecord
    sources_count INTEGER NOT NULL DEFAULT 0,
    sources_failed TEXT   NOT NULL DEFAULT '[]',
    elapsed_s     REAL    NOT NULL DEFAULT 0.0,
    created_at    TEXT    NOT NULL DEFAULT (datetime('now'))
);
\end{lstlisting}

\subsection{Caching policy}

Cache key: \texttt{SHA-256(source + ":" + canonicalize\_query(query))[:24]}.
TTL: 1 hour default (configurable via \texttt{CACHE\_TTL\_SECONDS}).
Backends: \texttt{InMemoryCache} (process-local, used in tests) and \texttt{FilesystemCache} (JSON files under \texttt{CACHE\_DIR}, survives restarts).

Cache hit rate in a demo run of 5 questions with 2 repeats: 40\% (2 of 5 questions hit the cache on the second run).

Staleness risk: Wikipedia summaries can change; the 1-hour TTL balances freshness vs.\ API politeness.

% =====================================================================
\section{Testing}\label{sec:testing}

\subsection{Strategy \& coverage}

Coverage: \textbf{$\geq$65\%} of \texttt{src/}.
All tests are fully offline — \texttt{ai.*} functions are monkey-patched via \texttt{conftest.py} fixtures (\texttt{FakeLLM}, \texttt{FakeAIService}, \texttt{FakeRepository}). HTTP is never called in the test suite.
Provided \texttt{tests/test\_ai\_smoke.py}: \textbf{passing}.

\begin{center}\small
\begin{tabular}{lc}
\toprule
\textbf{Suite} & \textbf{Coverage} \\
\midrule
\texttt{src/} (SE layer) & $\geq$65\,\% \\
Provided \texttt{ai/} smoke tests & passing \\
\bottomrule
\end{tabular}
\end{center}

\subsection{Notable tests}

\begin{enumerate}
  \item \textbf{Happy-path end-to-end} (\texttt{test\_researcher.py::test\_research\_returns\_result}): wires up \texttt{ResearchEngine} with \texttt{FakeAIService} and \texttt{FakeLLM}, calls \texttt{engine.research()}, asserts a \texttt{ResearchResult} with non-empty answer.
  \item \textbf{Concurrency} (\texttt{test\_concurrency.py::test\_one\_fails\_others\_succeed}): passes \texttt{[good(1), bad(), good(3)]} to \texttt{run\_concurrent}; asserts \texttt{results[1]} is a \texttt{ValueError} while \texttt{results[0]==1} and \texttt{results[2]==3}.
  \item \textbf{Error path} (\texttt{test\_researcher.py::test\_research\_partial\_source\_failure}): monkey-patches \texttt{fetch\_arxiv} to raise \texttt{ConnectionError}, asserts \texttt{"arxiv" in result.sources\_failed} and answer is non-empty.
\end{enumerate}

% =====================================================================
\section{Deployment \& Reproducibility}\label{sec:deploy}

\subsection{Docker image}

Multi-stage Dockerfile: builder stage installs deps into \texttt{/opt/venv}; runtime stage copies the venv and runs as non-root \texttt{appuser}.

\begin{verbatim}
docker build --platform linux/amd64 -t finalproj .
docker run finalproj                         # offline demo
docker run --env-file .env finalproj \
    python -m researcher ask "What is AI?"   # live CLI
\end{verbatim}

Image size: approximately 450 MB (python:3.12-slim base + dependencies).

\subsection{Configuration}

Every env var is listed in \texttt{.env.example}. Missing \texttt{LLM\_PROVIDER} defaults to \texttt{anthropic}; missing API key raises \texttt{ProviderError} on first synthesis call (not on startup — allows offline demo to work without any key).

% =====================================================================
\section{Limitations \& Future Work}\label{sec:limits}

\begin{itemize}
  \item \textbf{In-memory cache is lost on restart.} \texttt{FilesystemCache} survives restarts but uses plain JSON; a Redis-backed cache with atomic operations would be production-grade.
  \item \textbf{Web-search quota.} Tavily is the default web provider and requires \texttt{TAVILY\_API\_KEY}. If no key is provided, live web search fails clearly and the system degrades to the remaining sources.
  \item \textbf{SQLite is single-writer.} Horizontal scaling requires PostgreSQL (swap \texttt{DATABASE\_URL} to a \texttt{postgresql+asyncpg://} URL).
  \item \textbf{No token budget.} Very long questions may consume more LLM tokens than expected; a token-aware rate limiter would add cost control.
  \item \textbf{No persistent multi-provider failover.} If Anthropic is down, the system fails synthesis. A multi-provider fallback (Anthropic $\to$ OpenAI $\to$ Gemini) would improve resilience.
\end{itemize}

% =====================================================================
\section{Tools \& Acknowledgements}\label{sec:tools}

We used Cursor (Claude Sonnet 4.6) as follows:
\begin{itemize}
  \item \texttt{src/services/ai\_service.py}: Cursor scaffolded the tenacity retry decorator structure; we rewrote the per-call \texttt{asyncio.timeout()} integration after observing that outer timeouts do not isolate per-source failures.
  \item \texttt{src/concurrency/orchestrator.py}: Cursor helped review the concurrency pattern; we wired the production orchestrator through the semaphore-bounded \texttt{run\_concurrent} pipeline with per-source timeout isolation and graceful degradation.
  \item \texttt{tests/test\_concurrency.py}: Claude suggested test cases; we kept 5 of 6 and hand-wrote \texttt{test\_orchestrator\_cache\_hit}.
\end{itemize}

% =====================================================================
\section*{References}
\addcontentsline{toc}{section}{References}

\begin{enumerate}[label={[\arabic*]},leftmargin=2.4em,itemsep=2pt]
  \item Anthropic API documentation, \url{https://docs.anthropic.com/}
  \item tenacity library documentation, \url{https://tenacity.readthedocs.io/}
  \item pydantic-settings documentation, \url{https://docs.pydantic.dev/latest/concepts/pydantic_settings/}
  \item Wikipedia REST API, \url{https://en.wikipedia.org/api/rest_v1/}
  \item arXiv API documentation, \url{https://info.arxiv.org/help/api/index.html}
  \item Python asyncio documentation, \url{https://docs.python.org/3/library/asyncio.html}
  \item aiosqlite library, \url{https://aiosqlite.omnilib.dev/}
\end{enumerate}

\appendix
% =====================================================================
\section{Appendix --- Reproducing the benchmark}
\textbf{Exact commands to reproduce the \S\ref{sec:concurrency} benchmark.}

\begin{lstlisting}[style=python,language=bash,numbers=none]
git clone https://github.com/emiljafarov3841/async-research-assistant
cd async-research-assistant
cp .env.example .env  # fill in keys for live mode, or use --offline
pip install -r requirements.txt
python scripts/bench.py --n 5 --offline --delay 0.05
\end{lstlisting}

\vspace{1em}\hrule\vspace{0.6em}
\noindent\small\textit{End of report --- thank you for reading.}

\end{document}
